\barrerapagina
\section{Reducción dimensional para visualización}
La transformación PCA aprobada se ajustó con las 4\,599 filas de entrenamiento y se aplicó a las 2\,628 de validación, usando el mismo centrado. Las etiquetas solo colorean los puntos; no entran en el ajuste SVD. No se ajustó PCA con test ni se utilizaron las coordenadas bidimensionales para entrenar los clasificadores.

PC1 explica 32,987763\% y PC2 16,630823\%; juntas conservan 49,618586\%, por lo que 50,381414\% no se representa en el plano. La varianza total de entrenamiento es 20,0043497173 con $ddof=1$ y el rango numérico es 17. El rango menor que 20 es compatible con redundancias de las variables históricas; no impide conservar las dos primeras direcciones.

\begin{table}[htbp]
\centering\small
\caption{Cinco mayores coeficientes por magnitud en cada componente. El cuadrado es participación en la norma del vector, no importancia predictiva.}\label{tab:pca-cargas}
\begin{tabular}{lrrr}\toprule
Comp. & Característica & Coeficiente & Cuadrado (\%) \\ \midrule
PC1 & \texttt{\detokenize{unidades_media_3}} & 0.386556480 & 14.943 \\
PC1 & \texttt{\detokenize{unidades_mediana_3}} & 0.384643640 & 14.795 \\
PC1 & \texttt{\detokenize{unidades_lag_1}} & 0.379973488 & 14.438 \\
PC1 & \texttt{\detokenize{unidades_por_factura_lag_1}} & 0.377059706 & 14.217 \\
PC1 & \texttt{\detokenize{unidades_lag_2}} & 0.371108446 & 13.772 \\
PC2 & \texttt{\detokenize{facturas_lag_1}} & 0.494616346 & 24.465 \\
PC2 & \texttt{\detokenize{clientes_distintos_lag_1}} & 0.492434743 & 24.249 \\
PC2 & \texttt{\detokenize{variacion_relativa_1}} & 0.367933782 & 13.538 \\
PC2 & \texttt{\detokenize{cambio_unidades_1}} & 0.345813511 & 11.959 \\
PC2 & \texttt{\detokenize{tendencia_unidades_3}} & 0.300490423 & 9.029 \\
\bottomrule\end{tabular}
\end{table}

PC1 combina principalmente nivel histórico de unidades, con coeficientes positivos grandes en media, mediana, lags y unidades por factura. PC2 combina frecuencia comercial (facturas y clientes distintos) y cambios recientes de unidades. La convención de signo hace positivo el coeficiente de mayor magnitud: invertir un vector no modifica la varianza ni aporta una clase nueva. Los coeficientes negativos menores de PC2 incluyen unidades de los retardos segundo y tercero; su lectura es algebraica, no causal.

\figura{pca_train.png}{Proyección PCA de entrenamiento, coloreada por etiqueta real. Todos los puntos proceden del artefacto aprobado, sin selección ni desplazamiento artificial.}{fig:pca-train}{0.86}
\figura{pca_val.png}{Validación proyectada con media y componentes de entrenamiento. Se conserva la escala de ejes utilizada en la visualización aprobada.}{fig:pca-val}{0.86}
\figura{pca_varianza.png}{Varianza individual y acumulada de los componentes obtenidos de entrenamiento. El denominador incluye las 20 direcciones.}{fig:pca-varianza}{0.96}
Las clases se solapan visualmente. La pérdida de más de la mitad de la varianza limita cualquier inferencia sobre separabilidad en las 20 entradas. PCA no es un clasificador; el solapamiento no demuestra la causa del rendimiento ni invalida una señal fuera del plano. Ninguna línea bidimensional se presenta como frontera exacta de los modelos multidimensionales.
